\section{Mathematical preliminaries}
\label{sec:prelim}

This section fixes notation and the continuous-time language used for absorption models and birth--death processes later in the chapter. It is not a survey of probability: only those objects that reappear in subsequent sections are stated.

\subsection{Exponential clocks and competition}
\label{sec:exp}

\begin{definition}[Exponential law]
A non-negative random variable \(T\) is exponential with rate \(\lambda>0\), written \(T\sim\mathrm{Exp}(\lambda)\), if
\begin{equation*}
\pr{T>t}=\ee^{-\lambda t},\qquad t\ge0.
\end{equation*}
Then \(\mean{T}=1/\lambda\) and \(\mathrm{Var}(T)=1/\lambda^2\).
\end{definition}

The exponential law is memoryless: \(\pr{T>t+s\mid T>s}=\pr{T>t}\). Consequently, in continuous-time jump models, the waiting time until the next event can be redrawn from the current state without reference to the past.

\begin{proposition}[Competition of exponentials]
\label{prop:race}
Let \(T_1,\dots,T_m\) be independent, \(T_r\sim\mathrm{Exp}(a_r)\) with \(a_r>0\), and set \(\Lambda=\sum_{r=1}^m a_r\). Then
\begin{equation*}
\min_r T_r\sim\mathrm{Exp}(\Lambda),
\qquad
\pr{\text{event \(r\) occurs first}}=\frac{a_r}{\Lambda}.
\end{equation*}
\end{proposition}

In modelling language: each reaction channel carries an independent exponential clock; the first to ring determines the next jump, and the holding time is exponential with rate equal to the total propensity.

\subsection{Continuous-time Markov chains}
\label{sec:ctmc}

\begin{definition}[CTMC]
A continuous-time process \(\{\xx_t:t\ge0\}\) with discrete state space \(S\) is a continuous-time Markov chain if, for \(0\le t_0<\cdots<t_n<t_{n+1}\) and states \(i_0,\dots,i_{n+1}\in S\),
\begin{equation*}
\pr{\xx_{t_{n+1}}=i_{n+1}\mid \xx_{t_0}=i_0,\dots,\xx_{t_n}=i_n}
=\pr{\xx_{t_{n+1}}=i_{n+1}\mid \xx_{t_n}=i_n}.
\end{equation*}
\end{definition}

When transition probabilities depend only on the length of the time interval, the chain is time-homogeneous. The infinitesimal generator \(Q=(q_{ij})_{i,j\in S}\) is defined by
\begin{equation*}
q_{ij}=\lim_{\Delta t\to0^+}\frac{\pr{\xx_{\Delta t}=j\mid \xx_0=i}-\mathbf{1}_{\{i=j\}}}{\Delta t},
\qquad i,j\in S,
\end{equation*}
so that for \(i\neq j\), \(q_{ij}\ge0\) is the transition rate from \(i\) to \(j\), and \(q_{ii}=-\sum_{j\neq i}q_{ij}\). Holding times in state \(i\) are exponential with rate \(-q_{ii}\), and the next state is \(j\) with probability \(q_{ij}/(-q_{ii})\) when \(q_{ii}\neq0\).

\begin{definition}[Absorbing and transient states]
A state \(i\) is \emph{absorbing} if \(q_{ii}=0\). A non-absorbing state that is left almost surely and never returned to with positive probability is \emph{transient}. Extinction of a population, rupture of a host cell, and related terminal events are encoded as absorbing states in later models.
\end{definition}

Writing \(p_{ij}(t)=\pr{\xx_t=j\mid \xx_0=i}\), the Kolmogorov forward (master) equations read, in matrix form, \(\frac{\mathrm{d}}{\mathrm{d}t}P(t)=P(t)Q\), with \(P(0)=I\). Componentwise, for a countable state space,
\begin{equation}
\label{eq:master}
\ddt{p_j(t)} = \sum_{i\neq j} p_i(t)\,q_{ij} - p_j(t)\,(-q_{jj}),
\end{equation}
where \(p_j(t)=\pr{\xx_t=j}\) under a fixed initial law.

\subsection{Probability generating functions}
\label{sec:pgf}

\begin{definition}[Univariate PGF]
If \(X\) takes values in \(\{0,1,2,\dots\}\) with \(\pr{X=k}=p_k\), its probability generating function is
\begin{equation}
\label{eq:pgf-uni}
\phi(z)=\mean{z^X}=\sum_{k=0}^{\infty}p_k z^k,\qquad |z|\le1.
\end{equation}
\end{definition}

Factorial moments and probabilities are recovered by differentiation:
\begin{equation}
\label{eq:pgf-recover}
\mean{X}=\phi'(1),\qquad
\pr{X=k}=\frac{1}{k!}\,\phi^{(k)}(0),
\end{equation}
and \(\mathrm{Var}(X)=\phi''(1)+\phi'(1)-\bigl(\phi'(1)\bigr)^2\) when the derivatives exist at \(z=1\).

For a continuous-time process \(\{\xx_t\}\) on the non-negative integers, the time-dependent PGF
\begin{equation}
\label{eq:pgf-t}
G(z,t)=\mean{z^{\xt}}=\sum_{k\ge0}\pr{\xt=k}\,z^k
\end{equation}
converts the master equation~\eqref{eq:master} into a partial differential equation whenever the rates are polynomial in the state (as they are for linear birth--death and many catastrophe models). That conversion is used systematically in \cref{sec:moc}.

\begin{definition}[Multivariate and multi-type PGF]
\label{def:multi-pgf}
For a random vector \(\mathbf{X}=(X_1,\dots,X_d)\) with non-negative integer coordinates,
\begin{equation}
\label{eq:pgf-multi}
G(\mathbf{z})=\mean{z_1^{X_1}\cdots z_d^{X_d}}=\sum_{\mathbf{k}\in\mathbb{Z}_{\ge0}^d}\pr{\mathbf{X}=\mathbf{k}}\,z_1^{k_1}\cdots z_d^{k_d}.
\end{equation}
In a multi-type branching process, \(X_i\) is the count of type \(i\); the offspring law of each type is encoded by a vector of generating functions. Two-type models with conversion and type-dependent catastrophe rates appear in later chapters; here only the generating-function object is needed.
\end{definition}

\subsection{First-step analysis}
\label{sec:firststep}

Many extinction and absorption probabilities satisfy a closed equation obtained by conditioning on the first jump.

\begin{proposition}[First-step principle]
\label{prop:firststep}
Let \(\{\xx_t\}\) be a CTMC with absorbing set \(A\subset S\), and let \(h_i=\pr{\text{hit \(A\)}\mid \xx_0=i}\). Then \(h_i=1\) for \(i\in A\), and for transient \(i\),
\begin{equation}
\label{eq:firststep}
h_i = \sum_{j\neq i}\frac{q_{ij}}{-q_{ii}}\,h_j.
\end{equation}
The same idea in discrete time yields the Galton--Watson extinction iteration of \cref{sec:gw}: the extinction probability from one individual is the offspring PGF evaluated at that probability.
\end{proposition}

\subsection{Linear birth--death (pointer)}
\label{sec:bd-pointer}

The continuous-time linear birth--death process on \(\{0,1,2,\dots\}\) has rates
\begin{equation}
\label{eq:bd-rates}
q_{i,i+1}=i\lambda,\qquad q_{i,i-1}=i\mu,\qquad i\ge1,
\end{equation}
with \(0\) absorbing. Its extinction probability from one individual is \(\mu/\lambda\) when \(\lambda>\mu\) and \(1\) when \(\lambda\le\mu\). The associated PGF PDE and the quasi-stationary mean in the subcritical case are recorded in \cref{sec:cts}. Adding a global catastrophe (rupture) rate proportional to the load defines the birth--death--catastrophe models of Chapters~3--5; a discrete analogue is treated briefly in \cref{sec:bdc}.
